\begin{mistake}{2026-07-21}{01}

\begin{problemcontent}
设 \(f(x)\) 与 \(g(x)\) 在 \([a, b]\) 上连续，在 \((a, b)\) 内可导，且 \(f(a) = g(b) = 0\)，证明：至少存在一点 \(\xi \in (a, b)\)，使得
\[ f'(\xi) \int_{\xi}^b g(t) \mathrm{d}t + g'(\xi) \int_a^{\xi} f(t) \mathrm{d}t = 0 \]
\end{problemcontent}

\begin{solutioncontent}
构造辅助函数：
\[ F(x) = f(x) \int_x^b g(t) \mathrm{d}t + g(x) \int_a^x f(t) \mathrm{d}t \]
由于 \(f(x)\) 与 \(g(x)\) 在 \([a, b]\) 连续且在 \((a, b)\) 可导，所以 \(F(x)\) 在 \([a, b]\) 连续且在 \((a, b)\) 可导。

验证端点值：
\[ F(a) = f(a) \int_a^b g(t) \mathrm{d}t + g(a) \int_a^a f(t) \mathrm{d}t = 0 + 0 = 0 \]
\[ F(b) = f(b) \int_b^b g(t) \mathrm{d}t + g(b) \int_a^b f(t) \mathrm{d}t = 0 + 0 = 0 \]
故 \(F(a) = F(b) = 0\)。由罗尔定理可知，至少存在一点 \(\xi \in (a,b)\)，使得 \(F'(\xi) = 0\)。

对 \(F(x)\) 求导得：
\[ F'(x) = f'(x) \int_x^b g(t)\mathrm{d}t - f(x)g(x) + g'(x) \int_a^x f(t)\mathrm{d}t + g(x)f(x) \]
化简后：
\[ F'(x) = f'(x) \int_x^b g(t)\mathrm{d}t + g'(x) \int_a^x f(t)\mathrm{d}t \]
令 \(F'(\xi) = 0\)，即得所证等式。
\end{solutioncontent}

\mistakeknowledge{
\begin{itemize}
  \item \textbf{核心方法}：利用乘积求导公式凑微分，对各部分分别还原积函数，观察多余项能否相互抵消。
  \item \textbf{易错点}：对变限积分 \(\int_x^b g(t)\mathrm{d}t\) 求导时必须带出负号，漏掉将导致无法发现项的抵消。
  \item \textbf{关联知识}：罗尔定理；变上限与变下限积分求导法则。
\end{itemize}}

\end{mistake}
